\section{Contributions of the study}

In addressing the reasoning gap between large proprietary models and locally deployable Small Language Models in the Vietnamese legal domain, this thesis makes the following principal contributions:

\begin{enumerate}
    \item \textbf{A three-phase distillation and alignment pipeline for legal reasoning.} We propose and implement an end-to-end framework that combines Offline Logit Distillation, On-Policy Knowledge Distillation, and Diagnosis-Driven Preference Optimization across two independent tracks branching from a shared supervised fine-tuned base. The offline track transfers the teacher's knowledge and then performs targeted error correction via preference alignment on top of the offline checkpoint, while the on-policy track mitigates exposure bias through self-generated trajectories, yielding a compact 1.7B legal assistant that attains an Overall score of 0.5280 on the internal test set.

    \item \textbf{Empirical evidence that quantized adaptation prevents catastrophic forgetting.} We show that, on the 0.6B student, Quantized Low-Rank Adaptation (QLoRA) functions as an implicit regularizer that anchors prior supervised knowledge during distillation. This is evidenced by a steadily decreasing hard-label cross-entropy loss (settling around 0.37) under QLoRA, which stays consistently below the elevated loss (around 0.56) observed under standard LoRA, and a corresponding $\sim$33\% gain in syllogistic reasoning quality.

    \item \textbf{A diagnosis-driven preference-data construction protocol.} We replace conventional rule-based error scoring with a Large-Language-Model-as-a-Judge protocol that semantically evaluates both the reasoning trace and the final verdict. We demonstrate that this protocol produces substantially cleaner alignment data, increasing the rate of genuinely correct rollouts from 11.1\% to 20.7\% and improving the detection of substantively flawed reasoning from 40.5\% to 63.5\%, which directly strengthens the subsequent preference optimization.

    \item \textbf{A capacity-stratified empirical analysis across model scales.} Through a systematic comparison of the 0.6B and 1.7B students across every phase, we establish that model scale is the decisive factor for legal reasoning and, critically, that on-policy distillation and preference optimization are beneficial only above a minimum capacity threshold, destabilizing the sub-billion-parameter student. We further report an honest in-distribution versus out-of-distribution trade-off, showing that distillation improves in-distribution accuracy while the supervised baseline generalizes better on the external benchmark.

    \item \textbf{A reproducible evaluation on Vietnamese legal benchmarks.} We adapt and evaluate the distilled models on three core legal tasks (Natural Language Inference, Multiple-Choice, and Syllogism) across both an internal stratified test set and the external VLSP2025 Public-Test, providing a comprehensive characterization of distilled legal SLM behavior under both in-distribution and distribution-shift conditions.
\end{enumerate}
